\documentclass[11pt,a4paper]{article}
\usepackage[margin=0.85in,top=1in,bottom=1in]{geometry}
\usepackage{amsmath,amssymb,amsfonts}
\usepackage{graphicx}
\usepackage{float}
\usepackage{enumitem}
\usepackage{microtype}
\usepackage{caption}
\usepackage[hidelinks]{hyperref}
\usepackage[most,breakable]{tcolorbox}
\tcbuselibrary{breakable,skins}
\usepackage[utf8]{inputenc}
\usepackage[T1]{fontenc}
\usepackage{lmodern}
\captionsetup{font=small,labelfont=bf,skip=4pt}
\setlist{leftmargin=1.5em,topsep=2pt}

%% Breakable card — prevents content cutoff across pages
\newtcolorbox{solcard}[3]{
  breakable, enhanced,
  colback=white, colframe=gray!50, arc=2pt,
  title={\textbf{#1}\hfill\textsf{\small #3\,pts}},
  fonttitle=\normalsize\sffamily\bfseries,
  before upper={\small\textit{#2}\par\smallskip\hrule height 0.4pt\smallskip},
  top=4pt, bottom=6pt, left=7pt, right=7pt,
  parbox=false,
}

\title{\textbf{Double dumbbell --- Solution}}
\author{\normalsize X21 (2021) — English translation}
\date{}
\begin{document}\maketitle\vspace{-0.5em}
\begin{solcard}{A1}{Find the angular velocity vector of the second rod at the initial moment of time. Express your answer in terms of $\vec\omega$.}{0.30}

At the moment the rotation begins, the ball at the end of the second rod is motionless. Therefore, the rods begin to rotate in opposite directions. Since the balls are at rest and the lengths of the rods are the same $$\vec{\omega}_2=-\vec{\omega}_1 $$

\end{solcard}

\begin{solcard}{A2}{At the initial moment of time, find the angular velocity vector $\vec\omega_{rel_0}$ of the second rod in the reference system associated with the first rod. Express your answer using $\vec\omega$}{0.20}

From the law of addition of angular velocities $$\vec{\omega}_{rel}=\vec{\omega}_2-\vec{\omega}_1 $$ Combining the equality with the result of part $A1$ $$\vec{\omega}_{rel}=-2\vec{\omega} $$

\end{solcard}

\begin{solcard}{A3}{Find the potential energy of the load in the field of centrifugal force at a distance $r$ from the axis of rotation accurate to an arbitrary constant. Express your answer using $m$, $\omega$, $r$.}{0.50}

The potential energy can be obtained as the potential energy of a spring with a spring constant of $-m{\omega}^2$. Hence $$W_p=-\frac{m{\omega}^2r^2}{2}+C $$ In what follows, we will assume that the constant is equal to zero.

\end{solcard}

\begin{solcard}{A4}{Let the second rod rotate at an angle $\varphi$ relative to the first. Find its angular velocity $\omega_{rel}$ in this position. Express your answer using $\omega$ and $\varphi$.}{0.50}

As mentioned earlier, the Coriolis force does not change the kinetic energy of the system. Then the law of conservation of mechanical energy is satisfied in the system $$W_k+W_p=E $$ in initial moment time $$2m{\omega}^2L^2=W_k $$ $$-2m{\omega}^2L^2=W_p $$ Thus $$E=0 $$ in arbitrary moment time $$W_k=\frac{m{\omega_{rel}}^2L^2}{2} $$ $$W_p=-2m{\omega}^2L^2\cos^2\left(\frac{\varphi}{2}\right) $$ From which $$\vec{\omega}_{rel}=-2\vec{\omega}\cos\left(\frac{\varphi}{2}\right) $$

\end{solcard}

\begin{solcard}{A5}{Find the speed of the load in a fixed frame of reference. Express your answer in terms of $\omega$, $L$, $\varphi$.}{0.50}

One velocity component will be directed along the second rod, and the other will be perpendicular to it. We will receive $$v_1={\omega}L\sin(\varphi) $$ $$v_2={\omega}L(1+\cos(\varphi))-{\omega_{rel}}L $$ For square velocity obtain $$v^2={v_1}^2+{v_2}^2=(2\omega{L})^2\left(\left(\cos^2\left(\frac{\varphi}{2}\right)-\cos\left(\frac{\varphi}{2}\right)\right)^2+\sin^2\left(\frac{\varphi}{2}\right)\cos^2\left(\frac{\varphi}{2}\right)\right) $$ From which after simple transformations $$v=2{\omega}L\sqrt{2\left(\cos^2\left(\frac{\varphi}{2}\right)-\cos^3\left(\frac{\varphi}{2}\right)\right)} $$

\end{solcard}

\begin{solcard}{B1}{Let us denote the tension force of the rods as $T$ $$ $$ Find acceleration weight $a_{1l}$,$a_{1p}$ in projection on axis $OX$ and acceleration hinge $a_2$ in projection on axis $OY$.Answer express through $m$.$g$.$T$ and $\varphi$.}{0.40}

$$a_{1Lx}=\frac{T\cos(\varphi)}{m} $$ $$a_{1\perp x}=-\frac{T\cos(\varphi)}{m} $$ $$a_{2y}=-\left(\frac{T\sin(\varphi)}{m}+g\right) $$

\end{solcard}

\begin{solcard}{B2}{In the indicated position, find the speed of the weights $v_1$ and the speed of the hinge $v_2$.}{0.60}

From the law of conservation of mechanical energy $$\frac{2m}{2}({v_1}^2+{v_2}^2)=2m\left(\frac{{v_0}^2}{2}-gL\sin(\varphi)\right) $$ Equation kinematic relation $${v_1}={v_2}\tan(\varphi) $$ From which $$v_1=\sqrt{{v_0}^2-2gL\sin(\varphi)}\sin(\varphi) $$ $$v_2=\sqrt{{v_0}^2-2gL\sin(\varphi)}\cos(\varphi) $$

\end{solcard}

\begin{solcard}{B3}{Write down the equation of motion of the hinge relative to any of the masses, connecting $v_1$, $v_2$, $a_1$, $a_2$, $L$ and $\varphi$}{1.00}

The hinge moves in a circle relative to the loads. Hence $$\frac{v^2_{rel}}{L}={a_{1Lx}}\cos(\varphi)-{a_{2y}}\sin(\varphi) $$ Since velocity hinge and weight perpendicular $$v^2_{rel}={v_1}^2+{v_2}^2 $$ and finally $$\frac{{v_1}^2+{v_2}^2}{L}={a_{1Lx}}\cos(\varphi)-{a_{2y}}\sin(\varphi) $$

\end{solcard}

\begin{solcard}{B4}{Find the tension force $T$ from the equation obtained in step $B3$. Express the answer in terms of $m,{v_0},g,L,{\varphi}$}{0.70}

Combining the results of the first three points $$\left(\frac{{v_0}^2}{L}-2g\sin(\varphi)\right)=\frac{T}{m}+g\sin(\varphi) $$ from which $$T=m\left(\frac{{v_0}^2}{L}-3g\sin(\varphi)\right) $$

\end{solcard}

\begin{solcard}{B5}{Let us assume that at some moment the weights are lifted off the surface. $$ $$ Find value angle $\varphi_{neg}$ at this point in time.}{1.00}

Condition under which weights are lifted off the surface $$T\sin(\varphi)=mg $$ Substitute answer on part $B4$ $$m\left(\frac{{v_0}^2}{L}-3g\sin(\varphi)\right)\sin(\varphi)=mg $$ from which $$3\sin^2(\varphi)-\frac{{v_0}^2}{gL}\sin(\varphi)+1=0 $$ Find root square equation $$\sin(\varphi)=\frac{\frac{{v_0}^2}{gL}\pm\sqrt{\left(\frac{{v_0}^2}{gL}\right)^2-12}}{6} $$ First be achieved small root, and because final answer $$\sin(\varphi)=\frac{\frac{{v_0}^2}{gL}-\sqrt{\left(\frac{{v_0}^2}{gL}\right)^2-12}}{6} $$

\end{solcard}

\begin{solcard}{B6}{Find the minimum value $v_0$ at which the weights will come off the surface and the value of the angle $\varphi' $ at the moment of lifting}{0.50}

The discriminant of the quadratic equation must be non-negative, hence $$v_0\geq{\sqrt{2\sqrt{3}gL}} $$ The minimum is achieved in the case of zero of the discriminant, therefore $$sin(\varphi_{max})=\frac{1}{\sqrt{3}} $$

\end{solcard}

\begin{solcard}{C1}{At what values of $R$ does the collision of shells occur before their separation from the surface under the action of the rods? Express your answer using $L$,$\varphi_\text{ref}$,}{0.20}

Condition under which the collision occurs before the separation of $$R\geq{L}\cos(\varphi) $$

\end{solcard}

\begin{solcard}{C2}{Let immediately after the collision the projections of the speed of the load on the axes $OX$ and $OY$ is equal to to $v_{1x}$ and $v_{1y}$, respectively. $$ $$ Express projection velocity hinge on vertical axis $v_{2y}$ through $v_{1x}$,$v_{1y}$ and ${\varphi_0}$.}{1.00}

The load moves in a circle relative to the hinge, so $$\pm{v_{1x}}=(v_{1y}-v_{2y})\tan(\varphi_0) $$ $$v_{2y}=v_{1y}\pm{\frac{v_{1x}}{\tan(\varphi_0)}} $$

\end{solcard}

\begin{solcard}{C3}{Find the velocity of the hinge $v_2$ immediately after the impact. Express your answer using ${v_0},g,{L},R$}{1.30}

Let's solve the problem in the reference frame of the center of mass. Let $v_x$ and $v_y$ be the modulus of the relative speed of the loads and the modulus of the vertical velocity of the load relative to the hinge, respectively. Then the expression for kinetic energy in the reference frame of the center of mass is the following $$\frac{m{v_x}^2}{4}+\frac{m{v_y}^2}{2}=W_k $$ After the impact this also hold. $$ $$ From the equation kinematic relation $$\frac{v_x}{v_y}=tg(\varphi_0) $$ Immediately after impact this relation also hold. Therefore, from energy conservation follow, that after impact the magnitude of the relative velocity remains constant. This says that the weights exchange the horizontal velocity components with each other, and the hinge exchanges with they vertical constitute velocity. Thus, for impact hinge stops. $$v_{2}=0 $$

\end{solcard}

\begin{solcard}{C4}{Find the modulus of the load velocities $v_1$ immediately after the collision. Express the answer in terms of ${v_0},g,{L},R$}{1.30}

All kinetic energy goes to the loads. From the law of conservation of energy we get the answer $$v_{1}=\sqrt{{v_0}^2-2g\sqrt{L^2-R^2}} $$

\end{solcard}

\end{document}